\chapter{\texorpdfstring{$(P_5,\text{house})$}{(P5,house)}-free graphs without
\texorpdfstring{$C_5$}{C5}}
\label{cap:no_c5}

\section{Characterization of monopolarity}
\label{sec:no_c5_characterization}

There is a single additional minimal obstruction to monopolarity for $C_5$-free
graphs.

\begin{figure}[ht]
    \centering
    \begin{tikzpicture}[every node/.style={vertex}, node
        distance=1cm, on grid, edge]
        \foreach \i in {0,...,3}
            \node (\i)  at ({(360/4)*\i + 90}:1){};

            \node (4) at (0,0) {};
            \node (5) at ($(2) + (195:1)$) {};
            \node (6) at ($(2) + (255:1)$) {};
            \node (7) at ($(2) + (-15:1)$) {};
            \node (8) at ($(2) + (-75:1)$) {};
            \node [Label] (H) at (0,-2.5) {$H_{14}$};
            
            \foreach \i in {0,...,3}{
                \draw let \n1={int(mod(\i+1,4))} in (\i) to (\n1);
            }
            \draw (0) to (4);
            \draw (1) to (4);
            \draw (3) to (4);
            \draw (2) to (5);
            \draw (2) to (6);
            \draw (2) to (7);
            \draw (2) to (8);
            \draw (5) to (6);
            \draw (7) to (8);
    \end{tikzpicture}
    \caption{A minimal obstruction for monopolarity in $(P_5,\text{house})$-free
    graphs without an induced copy of $C_5$.}
    \label{fig:p5-house-obstruction-no-c5}
\end{figure}

\begin{proposition}
    \label{prop:minimal_obs_no_c5}
    The graph depicted in \Cref{fig:p5-house-obstruction-no-c5} is a minimal
    monopolar obstruction.
\end{proposition}

\begin{proof}
    Consider the following labelling of $H_{14}$ (see \Cref{fig:labelling_h14}).
    \begin{figure}[htb]
        \centering
        \begin{tikzpicture}[every node/.style={vertex}, node
            distance=1cm, on grid, edge, label distance=1.5pt]
            \node[label={above:$b$}] (0) at (90:1) {};
            \node[label={left:$c$}] (1) at (180:1) {};
            \node[label={[above right = 0pt and 2pt]15:$d$}] (2) at (270:1) {};
            \node[label={right:$a$}] (3) at (0:1) {};
            \node[label={below:$b'$}] (4) at (0,0) {};
            \node[label={left:$e$}] (5) at ($(2) + (195:1)$) {};
            \node[label={below:$f$}] (6) at ($(2) + (255:1)$) {};
            \node[label={right:$g$}] (7) at ($(2) + (-15:1)$) {};
            \node[label={below:$h$}] (8) at ($(2) + (-75:1)$) {};
            
            \foreach \i in {0,...,3}{
                \draw let \n1={int(mod(\i+1,4))} in (\i) to (\n1);
            }
            \draw (0) to (4);
            \draw (1) to (4);
            \draw (3) to (4);
            \draw (2) to (5);
            \draw (2) to (6);
            \draw (2) to (7);
            \draw (2) to (8);
            \draw (5) to (6);
            \draw (7) to (8);
        \end{tikzpicture}
        \caption{A labelling of $H_{14}$.}
        \label{fig:labelling_h14}
    \end{figure}
    Notice $d$ must be in $A$, otherwise, since $\l{e,f}$ and $\l{g,h}$ are
    anticomplete to each other, one would have that one of such sets is
    contained in $A$ for $B$ to be $P_3$-free, contradicting that $A$ is stable.
    Additionally, since $bb'$ is a chord edge of a diamond in $H_{14}$, it
    follows from \Cref{lem:monopolar_diamond} that $b$ and $b'$ must be in $B$.
    Since $d$ is in $A$ and $A$ is stable, it follows that $\p{a,b,c}$ is a
    $P_3$ in $B$, contradicting that $G[B]$ is a cluster. Therefore, $H_{14}$ is
    not monopolar, and in \Cref{fig:minimal_h14} the monopolar partitions of
    $H_{14} - v$ are shown.

    \begin{figure}[ht]
        \centering
        \begin{tikzpicture}
            \matrix[column sep=0.5cm]{
            \node{
                \begin{tikzpicture}[every node/.style={vertex}, node distance=1cm, on grid, edge]
                    \node (0) at ({90}:1){};
                    \node[Label] (1) at ({180}:1){};
                    \node (2) at ({270}:1){};
                    \node (3) at ({0}:1){};

                    \node (4) at (0,0) {};
                    \node (5) at ($(2)+(195:1)$) {};
                    \node (6) at ($(2)+(255:1)$) {};
                    \node (7) at ($(2)+(-15:1)$) {};
                    \node (8) at ($(2)+(-75:1)$) {};

                    \draw (0)--(3);
                    \draw (3)--(2);
                    \draw (0)--(4);
                    \draw (3)--(4);
                    \draw (2)--(5);
                    \draw (2)--(6);
                    \draw (2)--(7);
                    \draw (2)--(8);
                    \draw (5)--(6);
                    \draw (7)--(8);

                    \begin{scope}[on background layer, scale=1.5]
                        \draw[draw=none, fill=red!20] \convexpath{0,3}{3.5pt};
                        \draw[draw=none, fill=red!20] \convexpath{5,6}{3.5pt};
                        \draw[draw=none, fill=red!20] \convexpath{7,8}{3.5pt};
                        \node [fit=(2), fill=blue!20, inner sep=0.5pt, draw=none] {};
                        \node [fit=(4), fill=blue!20, inner sep=0.5pt, draw=none] {};
                    \end{scope}
                \end{tikzpicture}
                }; &
                \node{
                \begin{tikzpicture}[every node/.style={vertex}, node distance=1cm, on grid, edge]
                    \node (0) at ({90}:1){};
                    \node (1) at ({180}:1){};
                    \node[Label] (2) at ({270}:1){};
                    \node (3) at ({0}:1){};

                    \node (4) at (0,0) {};
                    \node (5) at ($(2)+(195:1)$) {};
                    \node (6) at ($(2)+(255:1)$) {};
                    \node (7) at ($(2)+(-15:1)$) {};
                    \node (8) at ($(2)+(-75:1)$) {};

                    \draw (0)--(1);
                    \draw (1)--(3);
                    \draw (3)--(0);
                    \draw (0)--(4);
                    \draw (1)--(4);
                    \draw (3)--(4);
                    \draw (5)--(6);
                    \draw (7)--(8);

                    \begin{scope}[on background layer, scale=1.5]
                        \draw[draw=none, fill=red!20] \convexpath{0,4}{3.5pt};
                        \node [fit=(6), fill=red!20, inner sep=0.5pt, draw=none] {};
                        \node [fit=(8), fill=red!20, inner sep=0.5pt, draw=none] {};
                        \node [fit=(1), fill=blue!20, inner sep=0.5pt, draw=none] {};
                        \node [fit=(3), fill=blue!20, inner sep=0.5pt, draw=none] {};
                        \node [fit=(5), fill=blue!20, inner sep=0.5pt, draw=none] {};
                        \node [fit=(7), fill=blue!20, inner sep=0.5pt, draw=none] {};
                    \end{scope}
                \end{tikzpicture}
                }; &
                \node{
                \begin{tikzpicture}[every node/.style={vertex}, node distance=1cm, on grid, edge]
                    \node (0) at ({90}:1){};
                    \node (1) at ({180}:1){};
                    \node (2) at ({270}:1){};
                    \node (3) at ({0}:1){};

                    \node[Label] (4) at (0,0) {};
                    \node (5) at ($(2)+(195:1)$) {};
                    \node (6) at ($(2)+(255:1)$) {};
                    \node (7) at ($(2)+(-15:1)$) {};
                    \node (8) at ($(2)+(-75:1)$) {};

                    \draw (0)--(1);
                    \draw (1)--(2);
                    \draw (2)--(3);
                    \draw (3)--(0);
                    \draw (2)--(5);
                    \draw (2)--(6);
                    \draw (2)--(7);
                    \draw (2)--(8);
                    \draw (5)--(6);
                    \draw (7)--(8);

                    \begin{scope}[on background layer, scale=1.5]
                        \draw[draw=none, fill=red!20] \convexpath{5,6}{3.5pt};
                        \draw[draw=none, fill=red!20] \convexpath{7,8}{3.5pt};
                        \node [fit=(1), fill=red!20, inner sep=0.5pt, draw=none] {};
                        \node [fit=(3), fill=red!20, inner sep=0.5pt, draw=none] {};
                        \node [fit=(0), fill=blue!20, inner sep=0.5pt, draw=none] {};
                        \node [fit=(2), fill=blue!20, inner sep=0.5pt, draw=none] {};
                    \end{scope}
                \end{tikzpicture}
                }; &
                \node{
                \begin{tikzpicture}[every node/.style={vertex}, node distance=1cm, on grid, edge]
                    \node (0) at ({90}:1){};
                    \node (1) at ({180}:1){};
                    \node (2) at ({270}:1){};
                    \node (3) at ({0}:1){};

                    \node (4) at (0,0) {};
                    \node[Label] (5) at ($(2)+(195:1)$) {};
                    \node (6) at ($(2)+(255:1)$) {};
                    \node (7) at ($(2)+(-15:1)$) {};
                    \node (8) at ($(2)+(-75:1)$) {};

                    \draw (0)--(1);
                    \draw (1)--(2);
                    \draw (2)--(3);
                    \draw (3)--(0);
                    \draw (0)--(4);
                    \draw (1)--(4);
                    \draw (3)--(4);
                    \draw (2)--(6);
                    \draw (2)--(7);
                    \draw (2)--(8);
                    \draw (7)--(8);

                    \begin{scope}[on background layer, scale=1.5]
                        \draw[draw=none, fill=red!20] \convexpath{0,4}{3.5pt};
                        \draw[draw=none, fill=red!20] \convexpath{2,7,8}{3.5pt};
                        \node [fit=(1), fill=blue!20, inner sep=0.5pt, draw=none] {};
                        \node [fit=(3), fill=blue!20, inner sep=0.5pt, draw=none] {};
                        \node [fit=(6), fill=blue!20, inner sep=0.5pt, draw=none] {};
                    \end{scope}
                \end{tikzpicture}
                }; \\
            };
        \end{tikzpicture}
        \caption{Monopolar partitions of $H_{14} - v$.}
        \label{fig:minimal_h14}
    \end{figure}

\end{proof}

Thus far, \textsc{Monopolarity} has been solved for those
$\p{P_5,\text{house}}$-free graphs that are either chordal or have an induced
copy of $C_5$. If a graph in this family is neither chordal nor it has an
induced copy of $C_5$, it follows that it must have an induced copy of $C_4$.
Thus, all graphs considered in this section, unless otherwise stated, are
assumed to have an induced copy of $C_4$, which will be denoted by
$C=\p{a,b,c,d,a}$, and $Z=\l{a,b,c,d}$. We start by proving the following lemmas
regarding the structure of the exact neighborhoods of $C$.

\begin{lemma}
    Let $G$ be a connected $(P_5,\text{house},C_5)$-free graph with and induced
    copy of $C_4$ and free of the graphs depicted in Figures
    \ref{fig:chordal-p5-free-obstructions},
    \ref{fig:p5-house-obstruction-with-c5} and
    \ref{fig:p5-house-obstruction-no-c5}. Let $X$ and $Y$ be different subsets
    of $Z$.
    \begin{lemenum}
        \item If $X$ consists of three vertices and $Y$ intersects it at
            non-consecutive vertices, then $N_X$ is anticomplete to $N_Y$.
            \label{lem:anti_1}
        \item If $X$ consists of two non-consecutive vertices and $Y$ consist of
            one vertex in $X$, then $N_X$ is anticomplete to $N_Y$.
            \label{lem:anti_2}
        \item If $X$ consists of three vertices and $Y$ consists of a vertex
            different from the middle vertex of $X$, then $N_X$ is anticomplete
            to $N_Y$. \label{lem:anti_3}
        \item If $X$ consists of one vertex and $Y$ consists of at most one
            vertex, consecutive to the one in $X$, then $N_X$ is anticomplete to
            $N_Y$. \label{lem:anti_4}
    \end{lemenum} 
\end{lemma}

\begin{proof}
    \leavevmode \par
    \begin{enumerate}[label=(\roman*)]
        \item Follows directly from \Cref{lem:c5_anticomp_4}.
        \item Suppose without loss of generality that $X=\l{a,c}$ and $Y=\l{c}$.
            If $x\in N_X$ is adjacent to $y\in N_Y$, then $\l{a,b,c,x,y}$
            induces a copy of the house. It follows that $N_X$ is anticomplete
            to $N_Y$.
        \item Suppose without loss of generality that $X=\l{a,b,c}$. If
            $Y=\l{a}$ or $Y=\l{c}$, as in the previous example, $x\in N_X$
            adjacent to $y\in N_Y$ implies that $\l{a,c,d,x,y}$ induces a copy
            of the house. If $Y=\l{d}$, then $x\in N_X$ adjacent to $y\in N_Y$
            implies that $\l{b,c,d,x,y}$ induces a copy of the house. Then,
            $N_{X}$ is anticomplete to $N_Y$.
        \item Suppose without loss of generality that $Y\subseteq\l{a}$ and
            $X=\l{b}$. If $x\in N_X$ is adjacent to $y\in N_Y$, then
            $\p{y,x,b,c,d}$ is a $P_5$. Therefore, $N_X$ is anticomplete to
            $N_Y$.
    \end{enumerate}
\end{proof}

\begin{lemma}
    Let $G$ be a connected $(P_5,\text{house},C_5)$-free graph with and induced
    copy of $C_4$ and free of the graphs depicted in Figures
    \ref{fig:chordal-p5-free-obstructions},
    \ref{fig:p5-house-obstruction-with-c5} and
    \ref{fig:p5-house-obstruction-no-c5}. Let $X$ and $Y$ be different subsets
    of $Z$.
    \begin{lemenum}
        \item If $X$ consists of three vertices and $Y$ consists of two
            non-consecutive vertices, one of them being the middle vertex of
            $X$, then $N_X$ is complete to $N_Y$. \label{lem:comp_1}
    \end{lemenum} 
\end{lemma}

\begin{proof}
    \leavevmode \par
    \begin{enumerate}[label=(\roman*)]
        \item Suppose without loss of generality that $X=\l{a,b,c}$ and
            $Y=\l{b,d}$. If $x\in N_X$ is not adjacent to $y\in N_Y$, then
            $\l{a,b,d,y,x}$ induces a copy of the house. Hence, $N_X$ is
            complete to $N_Y$.
    \end{enumerate}
\end{proof}

\begin{lemma}
    Let $G$ be a connected $(P_5,\text{house},C_5)$-free graph with and induced
    copy of $C_4$ and free of the graphs depicted in Figures
    \ref{fig:chordal-p5-free-obstructions},
    \ref{fig:p5-house-obstruction-with-c5} and
    \ref{fig:p5-house-obstruction-no-c5}. Let $X$ and $Y$ be different subsets
    of $Z$.
    \begin{lemenum}
        \item If $N_X$ consists of two consecutive vertices, then $N_X$ is
            empty. \label{lem:type_1}
        \item If $X=Z$, then $N_X$ is empty. \label{lem:type_2}
        \item If $X$ consists of three vertices, then $N_X$ is a clique.
            \label{lem:type_3}
        \item If $X$ consists of two non-consecutive vertices, then $G[N_X]$ is
            a cluster. \label{lem:type_4}
        \item Let $X$ and $Y$ consist of three vertices such that they intersect
            in two. Then, one between $N_X$ and $N_Y$ is empty.
            \label{lem:type_5}
        \item If $X$ consists of one vertex and $Y$ consists of one vertex
            non-consecutive to the one in $X$, then one between $N_X$ and $N_Y$
            is empty. \label{lem:type_6}
        \item If $X$ consists of one vertex, then $G[N_X]$ is either a cluster
            or split. \label{lem:type_7}
        \item Every connected component of $G[N_\varnothing]$ is homogeneous to
            vertices in $N_X$ with $X$ nonempty. Moreover, each connected
            component is anticomplete to all $N_X$ but one. \label{lem:type_8}
    \end{lemenum}
\end{lemma}

\begin{proof}
    \leavevmode \par
    \begin{enumerate}[label=(\roman*)]
        \item Suppose without loss of generality that $X=\l{a,b}$. If $x\in
            N_X$, then $\l{a,b,c,d,x}$ induces a copy of the house. Therefore,
            $N_X$ is empty.
        \item If $x\in N_X$, then $\l{a,b,c,d,x}$ induces a copy of $H_{11}$.
            Thus, $N_X$ is empty.
        \item Follows directly from \Cref{lem:c5_type_1}.
        \item Follows directly from \Cref{lem:c5_type_2}.
        \item Suppose without loss of generality that $X=\l{a,b,c}$ and
            $Y=\l{b,c,d}$. On the one hand, if $x\in N_X$ is adjacent to $y\in
            N_Y$, then $\l{a,b,c,d,x,y}$ induces a copy of $H_{13}$. On the
            other hand, if $x$ is not adjacent to $y$, then $\l{a,c,d,x,y}$
            induces a copy of the house. Therefore, one between $N_X$ and $N_Y$
            is empty.
        \item Suppose without loss of generality that $X=\l{a}$, so $Y=\l{c}$.
            If $x\in N_X$ is adjacent to $y\in N_Y$, then $\p{x,a,b,c,y,x}$ is a
            $C_5$. On the contrary, $x$ not adjacent to $y$ implies that
            $\p{x,a,b,c,y}$ is a $P_5$. Thus, one between $N_X$ and $N_Y$ is
            empty.

        \item Suppose without loss of generality that $X=\l{a}$. Since $G[N_X]$
            is not a cluster, let $x$, $y$ and $z$ be vertices in $N_{X}$ that
            induce a $P_3$, $(x,y,z)$. Since $G$ is $C_5$-free and
            $H_{11}$-free, it follows that $N_X$ is $(C_4,C_5)$-free.
            Nonetheless, $G[N_X]$ is not split, so there is an induced copy of
            $2K_2$ in $N_X$. 
            
            Now, for $u$ and $v$ adjacent vertices in $N_X$ different from $x$,
            $y$ and $z$, due to $H_4$, they cannot be anticomplete to the set
            $\l{x,y,z}$. Suppose without loss of generality that they are
            anticomplete to $\l{y,z}$ and $u$ is adjacent to $x$. Then,
            $\l{x,y,z,u,v,a}$ induces a copy of $H_{2}$ if $v$ is adjacent to
            $x$, while $\l{v,u,x,y,z}$ induces a copy of $P_5$ if $v$ is
            nonadjacent to $x$. Hence, $\l{u,v}$ is not anticomplete to
            $\l{y,z}$. Moreover, suppose $\l{u,v}$ is anticomplete to $y$. If
            $u$ is adjacent to $z$, then $\l{u,x,y,z,a}$ induces a copy of
            $H_{11}$. On the contrary, $y$ is adjacent to $z$, so
            $\p{u,x,y,z,v,u}$ is a $C_5$. This implies that at least one vertex
            of each $K_2$ (disjoint from $\p{x,y,z}$) is adjacent to $y$.
            Additionlly, an edge with $y$ as an endpoint cannot be part of a
            $2K_2$ in $N_X$. 
            
            Thus, let $u$ and $v$ be adjacent vertices of $N_X$, and suppose
            without loss of generality that $u$ is adjacent to $y$. Similarly,
            let $u'$ and $v'$ be adjacent vertices in $N_X$ such that $u'$ is
            adjacent to $y$, and $\l{u',v'}$ is anticomplete to $\l{u,v}$. If
            neither $v$ nor $v'$ are adajcent to $y$, then $\p{v,u,y,u',v'}$ is
            a $P_5$. If both $v$ and $v'$ are adjacent to $y$, then
            $\l{u,v,y,a,u',v'}$ induces a copy of $H_1$. Finally, if just one of
            $v$ and $v'$ is adjacent to $y$, then $\l{u,v,y,a,u',v'}$ induces a
            copy of $H_2$. It follows that $G[N_X]$ cannot be neither a cluster
            nor split.

        \item Since $G$ is connected, each connected component of
            $G[N_\varnothing]$ is not anticomplete to at least one exact
            neighborhood of $C$, $N_X$. Due to
            \Cref{lem:type_1,lem:type_2,lem:anti_4}, $X$ consists of either two
            consecutive vertices or three vertices. At any rate, vertices in
            $N_X$ are nonadajcent to at least one vertex of $C$, say $a$, and
            one between $b$ and $c$ is in $X$, say $b$. If adjacent vertices $u$
            and $v$ in $N_{\varnothing}$ are such that $v$ is adjacent to $x\in
            N_X$ and $u$ is nonadjacent to $x$, then $\p{u,v,x,b,a}$ is a $P_5$.
            Thus, every $K_2$ in $N_{\varnothing}$ is either complete or
            anticomplete to each vertex in an $N_X$. This shows that each
            component of $G[N_{\varnothing}]$ is homogeneous to vertices in
            $N_X$ for any $X$.

            Now suppose that $u\in N_{\varnothing}$ is adjacent to $b'\in
            N_{abc}$ and let $C_u$ be the connected component of $u$ in
            $G[N_\varnothing]$. Since $N_{abc}$ is nonempty, by
            \Cref{lem:type_7}, the only bags to which $N_\varnothing$ can be not
            anticomplete are $N_{acd}$, $N_{ac}$ and $N_{bd}$. On the one hand,
            if $C_u$ is not anticomplete to $N_{acd}$, then it must be that $u$
            is adjacent to a vertex $d'\in N_{acd}$. Thus, since $b'$ is
            nonadjacent to $d'$ due to \Cref{lem:anti_1}, $\l{a,b,b',d',u}$
            induces a copy of the house. Then, $C_u$ is anticomplete to
            $N_{acd}$. By the same reason, $C_u$ is anticomplete to $N_{ac}$. On
            the other hand, if $C_u$ is not anticomplete to $N_{bd}$, then there
            is an $x\in N_{bd}$ such that $u$ is adjacent to $x$. Since $x$ is
            adjacent to $b'$ due to \Cref{lem:comp_1}, it follows that
            $\l{u,b',x,c,d}$ is a house. Then, $C_u$ is anticomplete to
            $N_{bd}$. This shows that a connected component in $N_{\varnothing}$
            not anticomplete to a $3$-bag is anticomplete to all other bags.

            Finally, if $u\in N_{\varnothing}$ is adjacent to $x\in N_{ac}$ and
            $C_u$ is not anticomplete to $N_{bd}$, then $u$ must be adjacent to
            $y\in N_{bd}$. Thus, $\l{u,x,d,c,y}$ induces a copy of either a
            $C_5$ or the house, depending on whether $x$ and $y$ are adjacent or
            not. This shoes that a connected component in $N_{\varnothing}$ not
            anticomplete to a $2$-bag is anticomplete to all other bags. Thus,
            each connected component of $G[N_{\varnothing}]$ is anticomplete to
            all bags but one.
    \end{enumerate}
\end{proof}

As in the case of chordal graphs, it will be useful to give additional relations
between the exact neighborhoods provided that one 3-bag is nonempty, suppose
without loss of generality that $N_{abc}$.

\begin{lemma}
    Let $G$ be a connected $(P_5,\text{house},C_5)$-free graph with and induced
    copy of $C_4$ and free of the graphs depicted in Figures
    \ref{fig:chordal-p5-free-obstructions},
    \ref{fig:p5-house-obstruction-with-c5} and
    \ref{fig:p5-house-obstruction-no-c5}. Suppose that $N_{abc}$ is nonempty and
    let $b'\in N_{abc}$.
    \begin{lemenum}
        \item The set $N_{bd}$ is stable.
        \item The sets $G[N_a]$ and $G[N_c]$ are cluster.
    \end{lemenum} 
\end{lemma}

\begin{theorem}
    A $(P_5,\text{house},C_5)$-free graph $G$ with an induced copy of $C_4$ is
    monopolar if and only if it does not contain any of graphs depicted in
    Figures \ref{fig:chordal-p5-free-obstructions},
    \ref{fig:p5-house-obstruction-with-c5} and
    \ref{fig:p5-house-obstruction-no-c5}, as an induced subgraph.
\end{theorem}

\begin{proof}
    Since monopolarity is hereditary, it follows from
    \Cref{prop:minimal_obs_no_c5,prop:minimal_obs_c5,prop:minimal_obs_chordal}
    that $G$ does not contain the mentioned graphs as induced subgraphs.

    Due to \Cref{lem:type_1,lem:type_2}, the only possibly nonempty exact
    neighborhoods of the $4$-cycle $C$ are
    \[
        N_\varnothing, N_{a}, N_{b}, N_{c}, N_{d}, N_{ac}, N_{bd}, N_{abc}, N_{bcd},
        N_{acd} \text{ and } N_{abd}.
    \]
    We proceed by considering the following cases:

    \paragraph{Case 1: At least one 3-bag is nonempty.} Suppose without loss of
    generality that $N_{abc}$ is nonempty and let $b'\in N_{abc}$. By
    \Cref{lem:type_5}, both $N_{bcd}$ and $N_{abd}$ are empty, so the only other
    possibly nonempty 3-bag is $N_{acd}$. Additionally, \Cref{lem:anti_1}
    ensures that $N_{abc}$ is anticomplete to $N_{acd}$.

    If $x$ and $x'$ are adjacent vertices in $N_{bd}$, by \Cref{lem:comp_1}, $x$
    and $x'$ are adjacent to $b'$, so $\l{b,b',x,x',c,d}$ induces a copy of
    $H_{13}$. Therefore $N_{bd}$ is stable.

    By \Cref{lem:type_6}, suppose without loss of generality that $N_c$ is
    empty. If $N_a$ has vertices $x$, $x'$ and $x''$ that induce a $P_3$, then,
    by \Cref{lem:anti_3}, these three vertices are nonadjacent to $b'$, so
    $\l{x,x',x'',a,b,b'}$ induces a copy of $H_4$. Hence, $G[N_a]$ is a cluster.

    On the one hand, by \Cref{lem:anti_2}, $N_{ac}$ is anticomplete to $N_a$,
    while \Cref{lem:type_4} implies that $G[N_{ac}]$ is a cluster. On the other
    hand, $N_{abc}$ and $N_{acd}$ are cliques by \Cref{lem:type_3}, and these
    are anticomplete to $N_a$ and $N_{ac}$ by \Cref{lem:anti_3,lem:anti_1},
    respectively. Then, up to this point,
    \[
        \p{N_{bd} \cup \l{a,c} , N_{abc} \cup N_{acd} \cup N_{ac} \cup N_{a} 
        \cup \l{b,d}}
    \]
    is a partial monopolar partition of $G$.

    Once more, \Cref{lem:type_6} ensures that one between $N_{b}$ and $N_d$ is
    empty. We proceed by taking subcases:
    
    \paragraph{Case 1.1: Both $N_b$ and $N_d$ are empty.} Notice the
    aforementioned partition works in this case, with left to decide how to
    partition $N_{\varnothing}$. By \Cref{lem:type_8}, all connected components
    of $G[N_\varnothing]$ are anticomplete to all but one bag. Let $S'$ be the
    set of vertices of $G[N_{\varnothing}]$ that are in trivial components and
    are anticomplete to $N_{bd}$. Clearly $S'$ is stable and anticomplete to
    $N_{bd}$, so
    \[
        \p{N_{bd} \cup \l{a,c} \cup S', N_{abc} \cup N_{acd} \cup N_{ac} \cup N_{a} 
        \cup \l{b,d}}
    \]
    is a partial monopolar partition of $G$. Now let $C_u$ be the connected
    component of $u$ in $G[N_\varnothing\setminus S']$. Suppose $C_u$ is not
    anticomplete to a $3$-bag, and suppose without loss of generality that $C_u$
    is not anticomplete to $N_{abc}$. Let $u$ and $v$ be adjacent vertices of
    $C_u$. By \Cref{lem:type_8}, $C_u$ is homogeneous to $N_{abc}$, so both $u$
    and $v$ are adjacent to some $x$ in $N_{abc}$. Then, $\l{a,b,x,c,u,v}$
    induces a copy of $H_{4}$. This implies that all connected components in
    $G[N_\varnothing\setminus S']$, which are nontrivial, are anticomplete to
    all $3$-bags. Hence, vertices in $N_\varnothing \setminus S'$ are in
    connected components that are anticomplete to one $2$-bag and not
    anticomplete to the remaining $2$-bag.
    
    Suppose without loss of generality that $C_u$ is not anticomplete to
    $N_{ac}$. Let $u$ and $v$ be adjacent vertices of $C_u$. Since $C_u$ is
    homogeneous to $N_{ac}$, there exists a vertex $x$ in $N_{ac}$ such that $u$
    and $v$ are adjacent to $x$. If $y$ is a vertex in $N_{ac}$ adjacent to $x$,
    then, due to \Cref{lem:type_8}, either both $u$ and $v$ are adjacent to it,
    or none is adjacent to it. If the latter case happens, then
    $\l{a,c,x,y,u,v}$ induces a copy of $H_4$. Thus, both $u$ and $v$ are
    adjacent to $y$. It follows that $C_u$ is complete to the connected
    component of $x$ in $G[N_{ac}]$, $C_x$. Furthermore, if vertices in $C_u$
    are complete to another connected component of $G[N_{ac}]$, say the
    component of $x'\in N_{ac}$, then $\l{u,v,x,x',a,c,b,b'}$ induces a copy of
    $H_{12}$. Therefore, $C_u$ is anticomplete to all connected components of
    $G[N_{ac}]$ but one, to which is complete. Moreover, if $C_w$ is another
    connected of $G[N_\varnothing \setminus S']$ that is complete to $C_x$,
    then, since $C_w$ are nontrivial, it follows that there is $w$ and $w'$ in
    $C_w$ adjacent. Thus, $\l{a,b,b',c,x,u,v,w,w'}$ induces a copy of $H_{14}$.
    This implies that each component in $G[N_\varnothing \setminus S']$ not
    anticomplete to $N_{ac}$ is associated to a unique component in $G[N_{ac}]$,
    and this association is unique and such that the components are complete to
    each other. Notice each component $C_u$ in $G[N_\varnothing \setminus S']$
    is $C_5$-free since $G$ is, is $C_4$-free since $G$ is $H_{11}$-free, and is
    $2K_2$-free since $G$ is $H_{14}$-free. Then, each $C_u$ has a split
    partition into a stable set $S_u$ and a clique $K_u$. The clique, along with
    the connected component of $N_{ac}$ to which is complete, forms a clique,
    since $G[N_{ac}]$ is a cluster due to \Cref{lem:type_4}. Meanwhile, the
    stables $S_u$ are anticomplete to $N_{bd}$ since they were not anticomplete
    to $N_{ac}$.

    Finally, let $u$, $v$ and $w$ be in the same component of $G[N_\varnothing
    \setminus S']$ such that they are adjacent to vertex $x\in N_{bd}$ and they
    induce a $P_3$. Then, $\l{u,v,w,x,b,b'}$ induces a copy of $H_{4}$. Thus,
    all connected components of $G[N_\varnothing \setminus S']$ that are not
    anticomplete to $N_{bd}$ are clusters, and anticomplete to all other bags.
    At last, let $R$ be a subset of $N_\varnothing \setminus S'$ with one vertex
    per each connected component of $G[N_\varnothing \setminus S']$ that is not
    anticomplete to $N_{ac}$. Let $K'$ be the subset of $N_\varnothing\setminus
    S'$ that remains after removing all vertices of components $C_u$ with $u\in
    R$. It follows that
    \[
        \p{N_{bd} \cup \l{a,c} \cup S' \cup \bigcup_{u\in R} S_u, N_{abc} \cup 
        N_{acd} \cup N_{ac} \cup N_{a} \cup \l{b,d} \cup \bigcup_{u\in R} K_u
        \cup K'}
    \]
    is a monopolar partition of $G$. Notice this analysis on how to partition
    $N_\varnothing$ is independent of whether $N_b$ and $N_d$ are empty or not
    (even though we used the fact that $N_{abc}$ is nonempty). Thus, this
    partition for $N_{\varnothing}$ will be used throughout all subcases of case
    1.
    
    \paragraph{Case 1.2: $N_d$ is nonempty and $N_{acd}$ is empty.} Notice
    $G[N_d]$ is $C_5$-free since $G$ is, $C_4$-free due to $H_{11}$, and
    $2K_2$-free due to $H_{14}$. Then $G[N_d]$ is split, so there is a partition
    $(S,K)$ with $S$ stable and $K$ a clique, with the property that $S$ is
    maximal. By \Cref{lem:anti_2}, $N_d$ is anticomplete to $N_{bd}$, while
    \Cref{lem:anti_3} implies that $N_d$ is anticomplete to $N_{abc}$. Finally,
    suppose $K$ is not anticomplete to $N_{ac}$. Then, there exists $v\in K$ and
    $u\in N_d$ adjacent vertices and $v$ adjacent to $w\in N_{ac}$ (the
    existance of $u$ is due to the maximality of $S$). If $u$ is not adjacent to
    $w$, $\l{u,v,w,c,d}$ induces a copy of the house. If $u$ is not adjacent to
    $w$, then $\l{a,b,c,d,b',u,v,w}$ induces a copy of $H_{12}$. Then, it must
    be that $K$ is anticomplete to $N_{ac}$, so
    \[
        \p{N_{bd} \cup S \cup \l{a,c} \cup S' \cup \bigcup_{u\in R} S_u, N_{abc} 
        \cup K \cup N_{ac} \cup N_{a} \cup \l{b,d} \cup \bigcup_{u\in R} K_u \cup K'}
    \]
    is a monopolar partition of $G$.

    \paragraph{Case 1.3: $N_b$ is nonempty.} Observe first that the case for
    $N_d$ and $N_{acd}$ nonempty is analogous to the case where $N_b$ is
    nonempty, since we have already supposed that $N_{abc}$ is nonempty. Let $E$
    be the subset of vertices in $N_b$ that are complete to $N_{abc}$, and let
    $D= N_b \setminus E$. Suppose $u$ and $v$ are adjacent vertices in $D$.
    Since $G$ is $H_{4}$-free, it cannot be that none of $u$ and $v$ are
    adjacent to each vertex in $N_{abc}$. Nonetheless, since $u$ and $v$ are in
    $D$, it follows that there is a $u'$ in $N_{abc}$ that is adjacent to $u$
    and nonadjacent to $v$. Similarly, there is a $v'$ in $N_{abc}$ that is
    adjacent to $v$ and nonadjacent to $u$. Since $N_{abc}$ is a clique, it
    follows that $u'$ is adjacent to $v'$. Then, $\l{b,u,u',v,v'}$ induces a
    copy of $H_{11}$. Therefore, $D$ is stable. Additionally, if $u\in E$ is
    adjacent to $v\in D$, letting $v'$ to be a vertex in $N_{abc}$ that is
    nonadjacent to $v$ leads to $\p{v,u,v',c,d}$ being a $P_5$. Then, $E$ is
    anticomplete to $D$. Also, due to \Cref{lem:anti_2}, $D$ is anticomplete to
    $N_{bd}$, so
    \[
        \p{N_{bd} \cup \l{a,c} \cup D, N_{abc} \cup N_{acd} \cup N_{ac} \cup N_a 
        \cup \l{b,d}}
    \]
    is a partial monopolar partition of $G$. Let $F$ be the set of vertices in
    $E$ that are not anticomplete to $N_{ac}$, and let $I=E\setminus F$. Suppose
    $u$ and $v$ in $F$ are adjacent. Since they are in $E$, both are adjacent to
    $b'\in N_{abc}$. Since $u$ is in $F$, there is a $u'$ in $N_{ac}$ that is
    adjacent to $u$. Similarly, there is a $v'$ in $N_{ac}$ that is adjacent to
    $v$. If $u'=v'$, then $\l{b,b',u,v,u',c}$ induces a copy of $H_{13}$. Then,
    suppose $u'\neq v'$ and $u$ is nonadjacent to $v'$, as well as $v$ is
    nonadjacent to $u'$. Then, if $u'$ is nonadjacent to $v'$,
    $\p{a,u',u,v,v',a}$ is a $C_5$. Otherwise, $u'$ is adjacent to $v'$, so
    $\l{a,u',u,v,v'}$ induces a copy of the house. It follows that $F$ is
    stable. Furthermore,  if $u\in F$ is adjacent to $v\in I$, letting $u'$ be a
    vertex in $N_{ac}$ adjacent to $u$, we get that $\p{v,u,u',c,d}$ is a $P_5$.
    Therefore, $F$ is anticomplete to $I$. Since $E$ was anticomplete to $D$,
    $F$ is anticomplete to $D$, so
    \[
        \p{N_{bd} \cup \l{a,c} \cup D \cup F, N_{abc} \cup N_{acd} \cup N_{ac} 
        \cup N_a \cup \l{b,d}}
    \]
    is a partial monopolar partition of $G$. Finally, notice that $G[I]$ is
    split since $G$ is $C_5$-free, $I$ is $C_4$-free due to $H_{11}$, and is
    $2K_2$-free due to $H_{1}$. Then, there is a partition of $I$ into $S$ a
    stable set and $K$ a clique. Since $E$ is anticomplete to $D$ and $I$ is
    anticomplete to $F$, it follows that $S$ is anticomplete to both $D$ and
    $F$. Meanwhile, $K$ is anticomplete to $N_{abc}$, $N_{ac}$ and anticomplete
    to $N_a$ and $N_{cda}$ due to \Cref{lem:anti_4,lem:anti_3}, respectively. It
    follows that
    \begin{align*}
        \bigg( 
            & N_{bd} \cup \l{a,c} \cup D \cup F \cup S \cup S' \cup 
            \bigcup_{u\in R} S_u , \\ 
            & N_{abc} \cup N_{acd} \cup N_{ac} \cup N_a \cup \l{b,d} \cup K 
            \cup \bigcup_{u\in R} K_u \cup K' \bigg)
    \end{align*}
    is a monopolar partition of $G$.

    \paragraph{Case 2: All $3$-bags are nonempty.} Suppose $u$ and $v$ are
    adjacent vertices in $N_{ac}$. Then, $\p{a,u,c,d,a}$ is a $4$-cycle with a
    nonempty $3$-bag, since $v\in N_{auc}$, so one could proceed as in the
    previous case. Thus, consider that $N_{ac}$ is stable. Similarly, consider
    $N_{bd}$ is stable. As in the previous case, by \Cref{lem:type_6}, suppose
    without loss of generality that both $N_{c}$ and $N_{d}$ are empty.

    Recall that $N_a$ is anticomplete to $N_b$ by \Cref{lem:anti_4}. Suppose
    $x$, $y$ and $z$ in $N_a$ induce a $P_3$. Similarly, suppose, $x'$, $y'$ and
    $z'$ in $N_b$ induce a $P_3$. Then. $\l{a,x,y,z,b,x',y',z'}$ induces a copy
    of $H_{7}$. This implies that at most one between $G[N_a]$ and $G[N_b]$ is
    not a cluster. Suppose without loss of generality that $G[N_b]$ is a
    cluster. Similarly, if $\l{u,v,w,x}$ induces a copy of $2K_2$ in $N_a$ and
    $\l{u',v',w',x'}$ induces a copy of $2K_2$ in $N_b$, then
    $\l{a,u,v,w,x,b,u',v',w',x'}$ induces a copy of $H_{5}$. Since $N_a$ and
    $N_b$ are $(C_4,C_5)$-free, it follows that at most one between $G[N_a]$ and
    $G[N_b]$ is not split. Consider the following subcases:

    \paragraph{Case 2.1: $G[N_a]$ is split.} Let $(S,K)$ be a partition of $N_a$
    with $S$ a stable set and $K$ a maximal clique. By \Cref{lem:anti_2}, $N_a$
    is anticomplete to $N_{ac}$ and $N_b$ is anticomplete to $N_{bd}$. Then, $S$
    is anticomplete to $N_{ac}$. If $K$ is nonempty, due to its maximality, it
    has at least two vertices, $u$ and $v$. If just one of them is adjacent to
    $w\in N_{bd}$, say $u$, then $\l{u,v,a,b,w}$ induce a house. On the
    contrary, if both $u$ and $v$ are adjacent to $w$, then $\p{a,b,w,u,a}$ is a
    $4$-cycle with a nonempty $3$-bag, since $v\in N_{avw}$. Thus, it follows
    that $K$ is anticomplete to $N_{bd}$. At any rate,
    \[
        \p{N_{ac} \cup \l{b,d} \cup S, N_{bd} \cup N_b \cup \l{a,c} \cup K}
    \]
    is a partial monopolar partition of $G$, with $N_\varnothing$ still to be
    assessed. If $N_{\varnothing}$ is empty, then we are done. If not, by
    \Cref{lem:type_8}, each connected component of $N_\varnothing$ is not
    anticomplete to one, and only one, between $N_{ac}$ and $N_{bd}$. Suppose
    without loss of generality that $N_{\varnothing}$ is not anticomplete to
    $N_{ac}$, and let $D$ be the set of vertices in $N_\varnothing$ that are not
    anticomplete to $N_{ac}$. Let $u$ and $v$ be vertices in $N_{ac}$. Suppose
    there are nonadjacent vertices $u'$ and $v'$ in $D$ such that $u$ is
    adjacent to $u'$ and nonadjacent to $v'$, and $v$ is adjacent to $v'$ and
    nonadjacent to $u'$. Then, since $N_{ac}$ is stable, $\p{u,u',b,v,v'}$ is a
    $P_5$. This implies that the vertices in $N_{ac}$ are ordered by the
    connected components to which they are adjacent.
    
    Let $x$ be the vertex in $N_{ac}$ that is adjacent to all connected
    components of $D$. Notice that $x$ is complete to all such components, so
    $x$ is complete to $D$. Observe that $\p{a,b,c,x,a}$ is a $C_4$ with the set
    $D$ complete to $x$ and anticomplete to all other vertices of the $C_4$.
    Then, by \Cref{lem:type_7}, $G[D]$ is either a cluster or split. Let $u$ be
    in $D$. If $v$ and $w$ are adjacent vertices in $N_b$, then $x$ is adjacent
    to both, otherwise, $\p{u,x,c,b,v}$ or $\p{u,x,c,b,w}$ are a $P_5$. Then,
    $\p{a,b,v,x,a}$ is a $C_4$ with a nonempty 3-bag, since $w\in N_{bvx}$.
    Thus, $N_b$ is stable. Consider the following subcases:

    \paragraph{Case 2.1.1: $E = N_\varnothing\setminus D$ is empty.} Then $D =
    N_\varnothing$. If $G[D]$ is a cluster, then
    \[
        \p{N_{ac} \cup \l{b,d} \cup S, N_{bd} \cup N_b \cup \l{a,c} \cup K \cup D}
    \]
    is a monopolar partition of $G$. On the contrary, if $G[D]$ is not a
    cluster, then it is split. Additionally, due to $H_{7}$, $G[N_a]$ is a
    cluster. Let $(S',K')$ be a partition of $D$ with $S'$ a stable set and $K'$
    a maximal clique. If $K'$ is not trivial, then $x$ is the only vertex in
    $N_{ac}$ that is not anticomplete to $K'$, otherwise, one would get a $C_4$
    with a nonempty $3$-bag. Therefore,
    \[
        \p{N_{bd} \cup N_b \cup \l{a,c} \cup S' , N_{ac} \cup N_a \cup \l{b,d} 
        \cup K'}
    \]
    is a monopolar partition of $G$.

    \paragraph{Case 2.1.2: $E$ is nonempty.} Similar to the case with $D$, there
    is a vertex $y\in N_{bd}$ complete to $E$, $N_a$ is stable, and $G[E]$ is
    either a cluster or split. Moreover, due to $P_5$, $x$ is adjacent to $y$.
    Then, due to $H_{7}$ at most one between $G[D]$ and $G[E]$ is not a cluster.
    Similarly, due to $H_{5}$, at most one between $G[D]$ and $G[E]$ is not
    split. Suppose without loss of generality that $G[D]$ is split.

    If $G[E]$ is a cluster, then, as in the last partition, let $(S',K')$ be a
    split partition of $G[D]$ with $K'$ a maximal clique. Hence,
    \[
        \p{N_{bd} \cup N_b \cup \l{a,c} \cup S' , N_{ac} \cup N_a \cup \l{b,d} 
        \cup K' \cup E}
    \]
    is a monopolar partition of $G$. On the contrary, if $G[E]$ is not a
    cluster, then $G[D]$ is a cluster and $G[E]$ is split. Then, if $(S',K')$ is
    a split partition of $E$ with $K'$ a maximal clique, it follows that
    \[
        \p{N_{ac} \cup N_a \cup \l{b,d} \cup S' , N_{bd} \cup N_b \cup \l{a,c} 
        \cup K' \cup D}
    \]
    is a monopolar partition of $G$.

    \paragraph{Case 2.2: $G[N_a]$ is not split.} Then $G[N_b]$ is split. In
    particular, due to \Cref{lem:type_7}, it must be that $G[N_a]$ is a cluster.
    Then, $N_a$ and $N_b$ are swapped as in the previous case, so
    \[
        \p{N_{bd} \cup N_b \cup \l{a,c} \cup S , N_{ac} \cup \l{b,d} \cup K}
    \]
    is a partial monopolar partition of $G$, with $(S,K)$ a split partition of
    $N_b$ with $S$ a stable set and $K$ a maximal clique. All cases of
    $N_\varnothing$ explained in the previous subcase still give monopolar
    partitions just by swapping the stable set with the set inducing the
    cluster.
\end{proof}